\documentclass[11pt]{article}
\usepackage[a4paper,margin=1in]{geometry}
\usepackage{amsmath,amssymb,amsthm,mathtools}
\usepackage{bm}
\usepackage{booktabs}
\usepackage{array}
\usepackage{enumitem}
\usepackage[dvipsnames]{xcolor}
\usepackage{hyperref}
\usepackage{microtype}
\usepackage{comment}

\hypersetup{
  colorlinks=true,
  linkcolor=MidnightBlue,
  urlcolor=MidnightBlue
}

% Define \WITHSOLUTIONS before inputting this file to show solutions.
\newif\ifsolutions
\ifdefined\WITHSOLUTIONS
  \solutionstrue
\else
  \solutionsfalse
\fi

\newcommand{\R}{\mathbb{R}}
\newcommand{\Pbb}{\mathbb{P}}
\newcommand{\E}{\mathbb{E}}
\newcommand{\Cov}{\operatorname{Cov}}
\newcommand{\Var}{\operatorname{Var}}
\newcommand{\iid}{\mathrel{\overset{\mathrm{iid}}{\sim}}}
\newcommand{\ML}{\mathrm{ML}}
\newcommand{\rv}[1]{\widetilde{#1}}
\newcommand{\rvv}[1]{\widetilde{\bm{#1}}}
\newcommand{\given}{\,\middle|\,}

\ifsolutions
  \newenvironment{answer}{%
    \par\medskip\noindent\textcolor{MidnightBlue}{\textbf{Solution.}}\ }
    {\par\medskip}
  \newcommand{\workingspace}[1]{}
\else
  \excludecomment{answer}
  \newcommand{\workingspace}[1]{\vspace{#1}}
\fi

\newcommand{\keybox}[1]{%
  \begin{center}
  \fcolorbox{MidnightBlue}{MidnightBlue!5}{%
    \begin{minipage}{0.91\textwidth}#1\end{minipage}}
  \end{center}}

\title{\huge Recitation 3\\[0.35em]
  \normalsize Gaussian Random Vectors and Bayesian Models\\[0.35em]
  \large\ifsolutions Solutions and Teaching Notes\else Review and Practice Problems\fi}
\author{Xuanxi Zhang}
\date{September 25, 2026}

\begin{document}
\maketitle

\ifsolutions
\section*{Pacing guide (75 minutes)}
\begin{center}
\begin{tabular}{@{}lll@{}}
\toprule
Time & Activity & Suggested use \\
\midrule
0--25 min & Review & Gaussian vectors, mean MLE, Bayesian updating \\
25--37 min & Problem 1 & 6 min individual work, then discussion \\
37--55 min & Problem 2 & 9 min individual work, then discussion \\
55--69 min & Problem 3 & 7 min individual work, then discussion \\
69--75 min & Wrap-up or Problem 4 & Use Problem 4 only if the class is ahead \\
\bottomrule
\end{tabular}
\end{center}

\paragraph{Teaching emphasis.}
For Gaussian maximum likelihood, derive only the mean-vector estimate.  Treat
the covariance matrix as fixed during that derivation; the matrix derivative
for estimating the covariance is outside today's scope.  In the Bayesian
section, keep the parameter random and distinguish the prior, likelihood, and
posterior at every step.
\fi

\section{Gaussian random vectors}

A $d$-dimensional Gaussian random vector with mean $\bm{\mu}\in\R^d$ and
positive-definite covariance matrix $\bm{\Sigma}\in\R^{d\times d}$ is written
\[
  \rvv{x}\sim\mathcal N(\bm{\mu},\bm{\Sigma}).
\]
Its joint pdf is
\[
  f_{\bm{\mu},\bm{\Sigma}}(\bm{x})
  =\frac{1}{\sqrt{(2\pi)^d\det(\bm{\Sigma})}}
   \exp\!\left[-\frac12(\bm{x}-\bm{\mu})^{\mathsf T}
     \bm{\Sigma}^{-1}(\bm{x}-\bm{\mu})\right].
\]
The diagonal entry $\bm{\Sigma}_{jj}$ is the variance of coordinate $j$, and
$\bm{\Sigma}_{jk}$ is the covariance between coordinates $j$ and $k$.

Every linear transformation of a Gaussian random vector is Gaussian:
\[
  \bm{A}\rvv{x}+\bm{b}
  \sim\mathcal N\!\left(\bm{A}\bm{\mu}+\bm{b},
    \bm{A}\bm{\Sigma}\bm{A}^{\mathsf T}\right).
\]
\paragraph{Why?}
One characterization of a Gaussian random vector is that every linear
combination of its coordinates is a univariate Gaussian random variable.  Set
$\rvv{z}:=\bm{A}\rvv{x}+\bm{b}$.  For any fixed vector $\bm{c}$,
\[
  \bm{c}^{\mathsf T}\rvv{z}
  =(\bm{A}^{\mathsf T}\bm{c})^{\mathsf T}\rvv{x}
    +\bm{c}^{\mathsf T}\bm{b}.
\]
The right-hand side is a linear combination of the coordinates of
$\rvv{x}$ plus a constant, so it is Gaussian.  Since this holds for every
$\bm{c}$, $\rvv{z}$ is a Gaussian random vector.  Its parameters follow from
linearity of the mean and the covariance rules:
\begin{align*}
  \E[\rvv{z}]&=\bm{A}\E[\rvv{x}]+\bm{b}
    =\bm{A}\bm{\mu}+\bm{b},\\
  \Cov(\rvv{z})&=\bm{A}\Cov(\rvv{x})\bm{A}^{\mathsf T}
    =\bm{A}\bm{\Sigma}\bm{A}^{\mathsf T}.
\end{align*}

Consequently, every subvector is Gaussian.  If
\[
  \begin{bmatrix}\rvv{x}\\ \rvv{y}\end{bmatrix}
  \sim \mathcal N\!\left(
  \begin{bmatrix}\bm{\mu}_x\\ \bm{\mu}_y\end{bmatrix},
  \begin{bmatrix}
    \bm{\Sigma}_{xx} & \bm{\Sigma}_{xy}\\
    \bm{\Sigma}_{yx} & \bm{\Sigma}_{yy}
  \end{bmatrix}\right),
\]
then conditioning also preserves Gaussianity:
\begin{align*}
  \rvv{y}\mid\rvv{x}=\bm{x}
  \sim\mathcal N\!\big(&\bm{\mu}_y+
    \bm{\Sigma}_{yx}\bm{\Sigma}_{xx}^{-1}
      (\bm{x}-\bm{\mu}_x),\\[-0.2em]
    &\bm{\Sigma}_{yy}-
    \bm{\Sigma}_{yx}\bm{\Sigma}_{xx}^{-1}\bm{\Sigma}_{xy}\big).
\end{align*}

\paragraph{Why?}
Define
\[
  \bm{B}:=\bm{\Sigma}_{yx}\bm{\Sigma}_{xx}^{-1},
  \qquad
  \rvv{r}:=\rvv{y}-\bm{\mu}_y
     -\bm{B}(\rvv{x}-\bm{\mu}_x).
\]
The pair $(\rvv{x},\rvv{r})$ is jointly Gaussian because it is a linear
transformation of $(\rvv{x},\rvv{y})$.  Moreover,
\[
  \Cov(\rvv{r},\rvv{x})
  =\bm{\Sigma}_{yx}-\bm{B}\bm{\Sigma}_{xx}=\bm{0}.
\]
For jointly Gaussian vectors, zero covariance implies independence, so
$\rvv{r}$ is independent of $\rvv{x}$.  Its mean is zero and its covariance is
\[
  \Cov(\rvv{r})
  =\bm{\Sigma}_{yy}
   -\bm{\Sigma}_{yx}\bm{\Sigma}_{xx}^{-1}\bm{\Sigma}_{xy}.
\]
Finally, rewrite $\rvv{y}$ as
\[
  \rvv{y}=\bm{\mu}_y+\bm{B}(\rvv{x}-\bm{\mu}_x)+\rvv{r}.
\]
After conditioning on $\rvv{x}=\bm{x}$, the first two terms are fixed and the
independent Gaussian residual $\rvv{r}$ is unchanged.  This gives exactly the
conditional mean and covariance above.

\section{Maximum likelihood for a Gaussian mean vector}

Suppose
\[
  \rvv{x}_1,\ldots,\rvv{x}_n
  \iid\mathcal N(\bm{\mu},\bm{\Sigma}),
\]
where $\bm{\Sigma}$ is fixed and positive definite.  For the observed dataset
$X=\{\bm{x}_1,\ldots,\bm{x}_n\}$, the terms of the log-likelihood that depend
on $\bm{\mu}$ are
\[
  \ell(\bm{\mu})
  =C-\frac12\sum_{i=1}^n
  (\bm{x}_i-\bm{\mu})^{\mathsf T}\bm{\Sigma}^{-1}
  (\bm{x}_i-\bm{\mu}).
\]
The gradient and Hessian are
\[
  \nabla_{\bm{\mu}}\ell(\bm{\mu})
  =\bm{\Sigma}^{-1}\sum_{i=1}^n(\bm{x}_i-\bm{\mu}),
  \qquad
  \nabla_{\bm{\mu}}^2\ell(\bm{\mu})=-n\bm{\Sigma}^{-1}.
\]
Because $\bm{\Sigma}^{-1}$ is positive definite, the Hessian is negative
definite, so the unique maximizer is
\[
  \boxed{\bm{\mu}_{\ML}=\frac1n\sum_{i=1}^n\bm{x}_i.}
\]
Thus the ML estimate is obtained by averaging each coordinate.


\section{Bayesian parametric models}

In a Bayesian model, the parameter is itself random.  Let $\rv{\theta}$ be a
parameter and let $\rvv{x}$ represent the data.  The model specifies
\begin{align*}
  &\textbf{prior:} && f_{\rv{\theta}}(\theta),\\
  &\textbf{likelihood:} && f_{\rvv{x}\mid\rv{\theta}}(\bm{x}\mid\theta),\\
  &\textbf{posterior:} &&
  f_{\rv{\theta}\mid\rvv{x}}(\theta\mid\bm{x})
  =\frac{f_{\rvv{x}\mid\rv{\theta}}(\bm{x}\mid\theta)
          f_{\rv{\theta}}(\theta)}
         {f_{\rvv{x}}(\bm{x})}.
\end{align*}
The denominator is a normalizing constant, so as a function of $\theta$,
\[
  \boxed{\text{posterior}\ \propto\ \text{likelihood}\times\text{prior}.}
\]


Two conjugate updates will be useful today.  A prior is called
\textbf{conjugate} when the posterior belongs to the same distribution family
as the prior.

\paragraph{Gaussian prior and Gaussian likelihood.}
Suppose
\[
  \rv{\mu}\sim\mathcal N(m_0,\tau^2),
  \qquad
  \rv{y}\mid\rv{\mu}=\mu\sim\mathcal N(\mu,s^2),
\]
where $\tau^2$ describes prior uncertainty and $s^2$ describes sensor noise.
Bayes' rule gives, as a function of $\mu$,
\begin{align*}
  f_{\rv{\mu}\mid\rv{y}}(\mu\mid y)
  &\propto f_{\rv{y}\mid\rv{\mu}}(y\mid\mu)f_{\rv{\mu}}(\mu)\\
  &\propto
  \exp\!\left[-\frac{(y-\mu)^2}{2s^2}\right]
  \exp\!\left[-\frac{(\mu-m_0)^2}{2\tau^2}\right].
\end{align*}
Expanding the two squares and keeping the terms that involve $\mu$ gives
\[
  \log f_{\rv{\mu}\mid\rv{y}}(\mu\mid y)
  =-\frac12\left[
    \left(\frac1{\tau^2}+\frac1{s^2}\right)\mu^2
    -2\left(\frac{m_0}{\tau^2}+\frac{y}{s^2}\right)\mu
  \right]+C.
\]
Define
\[
  v_1:=\left(\frac1{\tau^2}+\frac1{s^2}\right)^{-1},
  \qquad
  m_1:=v_1\left(\frac{m_0}{\tau^2}+\frac{y}{s^2}\right).
\]
Completing the square turns the preceding expression into
\[
  -\frac{(\mu-m_1)^2}{2v_1}+C',
\]
which is the log-pdf of a Gaussian distribution.  Therefore
\[
  \boxed{\rv{\mu}\mid\rv{y}=y\sim\mathcal N(m_1,v_1).}
\]
The quantities $1/\tau^2$ and $1/s^2$ are called \textbf{precisions}.  The
posterior precision is their sum, and
\[
  m_1=\frac{s^2}{s^2+\tau^2}m_0
      +\frac{\tau^2}{s^2+\tau^2}y
\]
is a weighted average of the prior mean and the observation.  A more precise
source receives more weight.  Also $v_1<\tau^2$ and $v_1<s^2$, so combining
the prior with the observation reduces uncertainty.

\paragraph{Beta prior and Bernoulli likelihood.}
Suppose $\rv{\theta}\sim\operatorname{Beta}(a,b)$, so on $0<\theta<1$ its
prior density has the form
\[
  f_{\rv{\theta}}(\theta)\propto
  \theta^{a-1}(1-\theta)^{b-1}.
\]
Given $\rv{\theta}=\theta$, let the Bernoulli observations be conditionally
independent.  If the observed data contain $h$ heads and $t$ tails, their
likelihood is
\[
  p_{\rvv{x}\mid\rv{\theta}}(\bm{x}\mid\theta)
  =\prod_i\theta^{x_i}(1-\theta)^{1-x_i}
  =\theta^h(1-\theta)^t.
\]
Multiplying likelihood and prior gives
\begin{align*}
  f_{\rv{\theta}\mid\rvv{x}}(\theta\mid\bm{x})
  &\propto \theta^h(1-\theta)^t
    \theta^{a-1}(1-\theta)^{b-1}\\
  &=\theta^{a+h-1}(1-\theta)^{b+t-1}.
\end{align*}
This is the kernel of another beta density, so
\[
  \boxed{\rv{\theta}\mid\text{data}
    \sim\operatorname{Beta}(a+h,b+t).}
\]
The update simply adds the observed counts to the two beta parameters.  For
example, the uniform prior is $\operatorname{Beta}(1,1)$; after two heads it
becomes $\operatorname{Beta}(3,1)$.

\newpage
\section*{Practice problems}

\noindent Problems 2 and 3 are adapted from Exercises 5.15 and 6.14 of
\emph{Probability and Statistics for Data Science}.

\subsection*{Problem 1: MLE of a Gaussian mean vector \hfill [6 minutes]}
Let
\[
  \rvv{x}_1,\rvv{x}_2,\rvv{x}_3
  \iid\mathcal N(\bm{\mu},\bm{\Sigma}),
  \qquad
  \bm{\Sigma}=\begin{bmatrix}2&1\\1&2\end{bmatrix},
\]
and suppose the observed vectors are
\[
  \bm{x}_1=\begin{bmatrix}2\\1\end{bmatrix},\qquad
  \bm{x}_2=\begin{bmatrix}4\\0\end{bmatrix},\qquad
  \bm{x}_3=\begin{bmatrix}3\\2\end{bmatrix}.
\]
\begin{enumerate}[label=(\alph*),itemsep=0.4em]
  \item Write the log-likelihood as a function of $\bm{\mu}$, up to an
        additive constant.
  \item Differentiate the general $n$-observation expression and derive
        $\bm{\mu}_{\ML}$.
  \item Compute $\bm{\mu}_{\ML}$ for the observed dataset.
  \item Would replacing $\bm{\Sigma}$ by the identity matrix change this ML
        estimate? Briefly explain.
\end{enumerate}
\workingspace{4.8cm}

\begin{answer}
For these three observations,
\[
  \ell(\bm{\mu})
  =C-\frac12\sum_{i=1}^3
    (\bm{x}_i-\bm{\mu})^{\mathsf T}\bm{\Sigma}^{-1}
    (\bm{x}_i-\bm{\mu}).
\]
For general $n$,
\[
  \nabla_{\bm{\mu}}\ell(\bm{\mu})
  =\bm{\Sigma}^{-1}\sum_{i=1}^n(\bm{x}_i-\bm{\mu}).
\]
Since $\bm{\Sigma}$ is invertible, setting the gradient to zero gives
\[
  \sum_{i=1}^n\bm{x}_i-n\bm{\mu}=\bm{0}
  \quad\Longrightarrow\quad
  \bm{\mu}_{\ML}=\frac1n\sum_{i=1}^n\bm{x}_i.
\]
The negative-definite Hessian $-n\bm{\Sigma}^{-1}$ confirms that this is the
unique maximum.  For the observed data,
\[
  \bm{\mu}_{\ML}
  =\frac13\left(
    \begin{bmatrix}2\\1\end{bmatrix}
    +\begin{bmatrix}4\\0\end{bmatrix}
    +\begin{bmatrix}3\\2\end{bmatrix}\right)
  =\boxed{\begin{bmatrix}3\\1\end{bmatrix}}.
\]
Replacing the common covariance by the identity changes the likelihood
contours but not the maximizing mean.  The same factor
$\bm{\Sigma}^{-1}$ multiplies every residual and cancels from the score
equation.
\end{answer}

\newpage
\subsection*{Problem 2: Gaussian Bayesian model (Exercise 5.15) \hfill [9 minutes]}
A noisy sensor is used to estimate the temperature on top of a mountain.  Let
$\rv{\mu}$ denote the unknown temperature.  The prior and sensor model are
\[
  \rv{\mu}\sim\mathcal N(0,\sigma^2),
  \qquad
  \rv{y}\mid\rv{\mu}=\mu\sim\mathcal N(\mu,1).
\]
\begin{enumerate}[label=(\alph*),itemsep=0.45em]
  \item Show that $[\rv{\mu},\rv{y}]^{\mathsf T}$ is a Gaussian random vector
        with mean zero, and derive its covariance matrix.
  \item Derive the posterior distribution of $\rv{\mu}$ given $\rv{y}=y$.
  \item Let $y=-1$ and $\sigma^2=0.01$.  Describe or sketch the prior and
        posterior.  Does the prior express confidence that the temperature is
        close to $0$?  Does one reading change it much?
  \item Repeat part (c) for $y=-1$ and $\sigma^2=10$.
\end{enumerate}
\workingspace{6.2cm}

\begin{answer}
Write $\rv{y}=\rv{\mu}+\rv{\varepsilon}$, where
$\rv{\varepsilon}\sim\mathcal N(0,1)$ is independent of $\rv{\mu}$.  Then
\[
  \begin{bmatrix}\rv{\mu}\\\rv{y}\end{bmatrix}
  =\begin{bmatrix}1&0\\1&1\end{bmatrix}
   \begin{bmatrix}\rv{\mu}\\\rv{\varepsilon}\end{bmatrix},
\]
so it is a linear transformation of a Gaussian random vector and is therefore
Gaussian.  Its mean is zero and its covariance matrix is
\[
  \begin{bmatrix}
    \Var(\rv{\mu}) & \Cov(\rv{\mu},\rv{y})\\
    \Cov(\rv{y},\rv{\mu}) & \Var(\rv{y})
  \end{bmatrix}
  =\boxed{\begin{bmatrix}
    \sigma^2 & \sigma^2\\
    \sigma^2 & \sigma^2+1
  \end{bmatrix}}.
\]
Using the conditional-Gaussian formula,
\[
  \boxed{
  \rv{\mu}\mid\rv{y}=y
  \sim\mathcal N\!\left(
    \frac{\sigma^2}{\sigma^2+1}y,
    \frac{\sigma^2}{\sigma^2+1}\right).}
\]
The posterior mean is the sensor reading shrunk toward the prior mean $0$.

For $y=-1$ and $\sigma^2=0.01$, the prior is
$\mathcal N(0,0.01)$ and the posterior is
\[
  \mathcal N\!\left(-\frac1{101},\frac1{101}\right)
  \approx\mathcal N(-0.0099,0.00990).
\]
The prior standard deviation is $0.1$, and the posterior standard deviation is
about $0.0995$.  The two narrow curves almost overlap.  This prior expresses
strong confidence near $0$, so one noisy reading barely changes it.

For $y=-1$ and $\sigma^2=10$, the prior is $\mathcal N(0,10)$ and the posterior
is
\[
  \mathcal N\!\left(-\frac{10}{11},\frac{10}{11}\right)
  \approx\mathcal N(-0.909,0.909).
\]
The prior standard deviation is about $3.16$, so it does not express strong
confidence near $0$.  The posterior standard deviation is about $0.953$ and
its center is close to the reading, so the posterior is much more concentrated
and meaningfully different from the prior.
\end{answer}

\newpage
\subsection*{Problem 3: Two coin flips (Exercise 6.14) \hfill [7 minutes]}
Let the probability of heads be a random variable
$\rv{\theta}\sim\operatorname{Uniform}[0,1]$.  Given $\rv{\theta}=\theta$,
the two flips $\rv{x}_1$ and $\rv{x}_2$ are conditionally independent
Bernoulli random variables with parameter $\theta$, where $1$ denotes heads.
\begin{enumerate}[label=(\alph*),itemsep=0.45em]
  \item Compute the joint pmf of $\rv{x}_1$ and $\rv{x}_2$.
  \item Both flips are observed to be heads.  Compute
        $\Pbb(\rv{\theta}<1/2\mid\rv{x}_1=1,\rv{x}_2=1)$.
  \item Are the two flips independent before conditioning on $\rv{\theta}$?
        Use your answer to part (a) to justify your conclusion.
\end{enumerate}
\workingspace{5.4cm}

\begin{answer}
For $a,b\in\{0,1\}$, integrate out the shared random parameter:
\[
  p_{\rv{x}_1,\rv{x}_2}(a,b)
  =\int_0^1 \theta^{a+b}(1-\theta)^{2-a-b}\,d\theta.
\]
Therefore
\[
\begin{array}{c|cc}
 & \rv{x}_2=0 & \rv{x}_2=1\\
\hline
\rv{x}_1=0 & \frac13 & \frac16\\[0.2em]
\rv{x}_1=1 & \frac16 & \frac13
\end{array}
\]
The posterior density after two heads is proportional to $\theta^2$ on
$[0,1]$.  Since $\int_0^1\theta^2\,d\theta=1/3$,
\[
  f_{\rv{\theta}\mid\rv{x}_1,\rv{x}_2}(\theta\mid1,1)=3\theta^2,
  \qquad 0\leq\theta\leq1.
\]
Hence
\[
  \Pbb(\rv{\theta}<1/2\mid\rv{x}_1=1,\rv{x}_2=1)
  =\int_0^{1/2}3\theta^2\,d\theta
  =\boxed{\frac18}.
\]
Each flip has marginal probability $1/2$ of heads, but
\[
  \Pbb(\rv{x}_1=1,\rv{x}_2=1)=\frac13
  \neq \frac12\cdot\frac12.
\]
Thus the flips are not marginally independent.  They are independent only
after conditioning on their shared parameter $\rv{\theta}$.
\end{answer}

\newpage
\subsection*{Problem 4 (backup): True or false?}
\noindent\textit{Optional if time allows. Suggested working time: 5--7 minutes.}

\begin{enumerate}[label=(\alph*),itemsep=0.65em]
  \item For a jointly Gaussian vector, two coordinates with covariance zero
        are independent.
  \item If every observation has the same known covariance matrix, changing
        that matrix changes the ML estimate of the Gaussian mean vector.
  \item In a Bayesian model, the posterior is proportional to the prior divided
        by the likelihood.
  \item Random variables may be conditionally independent given a shared
        parameter but dependent after that parameter is integrated out.
\end{enumerate}
\workingspace{4.2cm}

\begin{answer}
The answers are \textbf{true, false, false, true}.
\begin{enumerate}[label=(\alph*),itemsep=0.35em]
  \item This implication is a special property of jointly Gaussian random
        vectors; it is not true for arbitrary distributions.
  \item With a common positive-definite covariance, the score equation reduces
        to $\sum_i(\bm{x}_i-\bm{\mu})=\bm{0}$, so the estimate remains the
        coordinatewise sample mean.
  \item The posterior is proportional to likelihood times prior.
  \item Exercise 6.14 is an example: the flips are conditionally independent
        given $\rv{\theta}$ but positively dependent marginally.
\end{enumerate}
\end{answer}

\end{document}
